%% ma: L11-B15-0092
%% lop: 11
%% chuong: 5
%% bai: 15
%% dang: 11.15.5
%% loai: TLN
%% muc: VD
%% dap_an: "3"
%% nguon: "tai-lieu-hoc-tap-mon-toan-lop-11-theo-chuong-trinh-giao-duc-pho-thong-2018.pdf (Lớp Toán Thầy Duy, Phạm Lê Duy, THPT Chu Văn An), Chương 3 Bài 1, Phần 3, Câu 14 (tr.343)"
%% kien_thuc: "Lớp 11 Bài 15 (giới hạn của dãy số)"
%% trang_thai: cho-duyet
%% ly_do: "Dạng mới đề xuất (11.15.x), chờ giáo viên duyệt dạng; tài liệu gốc không có lời giải mẫu, đáp án do kho tự giải và kiểm bằng sympy."
\begin{ex}
Biết rằng giới hạn \[\lim\left(1+3n-\sqrt{9n^2-n+7}\right)=\lim\left(a+\dfrac{n-7}{b\cdot n+c\cdot\sqrt{9n^2-n+7}}\right)\] với $a;b;c$ là các số tự nhiên. Tính giá trị $|a-b|+c$.
\shortans{3}
\loigiai{\sloppy \[3n-\sqrt{9n^2-n+7}=\dfrac{9n^2-(9n^2-n+7)}{3n+\sqrt{9n^2-n+7}}=\dfrac{n-7}{3n+\sqrt{9n^2-n+7}}.\] Vậy $a=1$, $b=3$, $c=1$ và $|a-b|+c=2+1=3$.}
\end{ex}
